\documentclass[12pt,a4paper]{article}

% Поддержка русского текста для любого доступного движка LaTeX
\usepackage{iftex}
\ifPDFTeX
    \usepackage[T2A]{fontenc}
    \usepackage[utf8]{inputenc}
    \usepackage[russian]{babel}
\else
    \usepackage{fontspec}
    \usepackage{polyglossia}
    \setmainlanguage{russian}
\fi

\usepackage[
    a4paper,
    margin=2.5cm
]{geometry}

\usepackage{listings}
\usepackage{float}

\lstset{
    language=C,
    basicstyle=\ttfamily\small,
    numbers=left,
    numberstyle=\tiny,
    frame=single,
    breaklines=true,
    showstringspaces=false
}

\title{Документация к GPU Sandbox}

\author{}

\date{}

\begin{document}

\maketitle

\tableofcontents

\newpage

\section{Назначение программы}

GPU Sandbox представляет собой программную модель вычислительного процессора.

\section{Архитектура}

На текущем этапе модель содержит одно вычислительное ядро,
память и программу, состоящую из последовательности инструкций.

\subsection{Memory}

\texttt{data[MEMORY\_SIZE]} --- это память виртуального GPU.
Именно в ней хранятся все данные, с которыми выполняются операции.

Количество элементов памяти определяется переменной
\texttt{MEMORY\_SIZE}. Изначально её значение равно 1 МБ.
В будущем это значение будет изменено.

\subsection{Core}

\texttt{Core} --- ядро процессора. В нём находятся следующие компоненты.

\subsubsection{Регистры}

\texttt{registers[REGISTERS\_COUNT]} --- это массив регистров,
которые хранят в себе промежуточную информацию, с которой
впоследствии будут произведены операции.

Количество регистров определяется переменной
\texttt{REGISTERS\_COUNT}, которая изначально равна 8.
В будущем это значение будет изменено.

\subsubsection{Программный счётчик}

\texttt{pc} --- программный счётчик (от англ. \textit{Program Counter}).
Это специальный компонент, который указывает на инструкцию
(команду), которую видеокарта должна выполнить следующей.

Максимальное значение счётчика команд измеряется длиной текущей
инструкции.

\subsubsection{Текущее состояние}

\texttt{current\_state} --- это компонент, отвечающий за текущее
состояние ядра.

Состояния бывают следующие:

\begin{enumerate}

    \item \texttt{NORMAL} --- нормальное состояние ядра.
    Это состояние, которое принимает ядро во время работы при
    условии, что всё работает в штатном режиме без ошибок.

    \item \texttt{JUMP} --- состояние, сигнализирующее о том,
    что ядро перепрыгивает с инструкции, которую оно должно было
    исполнить следующей, на инструкцию с индексом $n$.
    
    Подробнее данная операция рассматривается ниже.

    \item \texttt{HALT} --- состояние, сигнализирующее
    о прекращении выполнения программы ядром.

    \item \texttt{ERROR} --- состояние, которое означает,
    что ядро столкнулось с какой-либо ошибкой, из-за чего
    экстренно завершило свой процесс.
    
    Обычно появляется при неправильно заданной инструкции.

    \item \texttt{TIMEOUT} --- состояние, схожее с
    \texttt{ERROR}, но отличается тем, что программа превысила
    максимально возможное количество циклов.

    Один цикл равен одной инструкции. Максимальное количество
    циклов до перехода ядра в состояние \texttt{TIMEOUT}
    определяется переменной \texttt{MAX\_CYCLES}.
    
    Изначально её значение равно 1\,000\,000.
    В будущем это значение может быть изменено.

\end{enumerate}

\subsection{Instructions}

\texttt{Instructions} --- это инструкции (команды), которые
должно выполнить ядро.

Для грамотной работы с инструкциями также нужны следующие
компоненты.

\subsubsection{Операнды}

\texttt{a}, \texttt{b}, \texttt{c} --- операнды, в которых
записывается информация, которую будет использовать ядро
для выполнения текущей операции. Данные наименования операндов в будущем могут быть изменены

\subsubsection{Opcode}

\texttt{Opcode} (образовано от \textit{Operation Code}) ---
это номер операции, который определяет, какую операцию нужно выполнить. Наглядная визуализация распределения операндов для работы опокода находится в таблице 1.

На данный момент в программе присутствуют следующие опкоды:

\begin{enumerate}

    \item \texttt{OP\_LOAD} --- операция, которая позволяет
    считать из памяти значение по адресу, хранящемуся в регистре \texttt{b},
    и записать его в регистр \texttt{a}.

    \item \texttt{OP\_STORE} --- операция, позволяющая записать
    значение регистра \texttt{b} в память по адресу, хранящемуся
    в регистре \texttt{a}.

    \item \texttt{OP\_JUMP} --- операция, позволяющая перейти
    на индекс инструкции, сохранённый в регистре \texttt{a}.

    \item \texttt{OP\_ADD} --- операция, которая записывает
    в регистр \texttt{a} сумму значений из регистров
    \texttt{b} и \texttt{c}.

    \item \texttt{OP\_SUB} --- операция, которая вычитает
    значение регистра \texttt{c} из регистра \texttt{b}
    и записывает результат в регистр \texttt{a}.

    \item \texttt{OP\_HALT} --- операция, означающая конец
    выполнения программы. Она переводит ядро в состояние
    \texttt{HALT}.
\end{enumerate}

    \begin{table}[H]
    \centering
    \caption{Набор инструкций виртуального процессора}
    \label{tab:opcodes}
    
    \begin{tabular}{|c|c|c|c|p{6cm}|}
        \hline
        \textbf{Opcode} & \textbf{a} & \textbf{b} & \textbf{c} & \textbf{Назначение} \\
        \hline
        
        \texttt{OP\_LOAD}
        & регистр
        & регистр с адресом
        & ---
        & Загрузка значения из памяти по адресу из \texttt{b} в \texttt{a}. \\
        \hline
        
        \texttt{OP\_STORE}
        & регистр с адресом
        & регистр со значением
        & ---
        & Запись значения из \texttt{b} в память по адресу из \texttt{a}. \\
        \hline
        
        \texttt{OP\_JUMP}
        & регистр с индексом инструкции
        & ---
        & ---
        & Переход к инструкции, индекс которой хранится в \texttt{a}. \\
        \hline
        
        \texttt{OP\_ADD}
        & регистр
        & регистр
        & регистр
        & Сложение значений регистров \texttt{b} и \texttt{c} с записью результата в \texttt{a}. \\
        \hline
        
        \texttt{OP\_SUB}
        & регистр
        & регистр
        & регистр
        & Вычитание значения регистра \texttt{c} из регистра \texttt{b} с записью результата в \texttt{a}. \\
        \hline
        
        \texttt{OP\_HALT}
        & ---
        & ---
        & ---
        & Прекращение выполнения программы и переход ядра в состояние \texttt{HALT}. \\
        \hline
        
    \end{tabular}
\end{table}



\end{document}
